\subsection{Complete filtered modules}

PROPOSITION 5. Let G be a filtered group whose filtration (G \(_{n}\) ) consists of invariant subgroups of G. The following conditions are equivalent:

\begin{enumerate}
\def\labelenumi{(\alph{enumi})}
\item
  G is a complete topological group.
\item
  The associated Hausdorff group \(\mathbf{G}' = \mathbf{G} / \left(\bigcap_{n\in \mathbf{Z}}\mathbf{G}_n\right)\) is complete.
\item
  Every Cauchy sequence in G is convergent.
\end{enumerate}

If \(\mathbf{G}\) is commutative and written additively, these conditions are also equivalent to the following:

\begin{enumerate}
\def\labelenumi{(\alph{enumi})}
\setcounter{enumi}{3}
\tightlist
\item
  Every family \((x_{\lambda})_{\lambda \in L}\) of elements of \(\mathbf{G}'\) which converges to 0 with respect to the filter \(\mathfrak{F}\) of complements of finite subsets of \(\mathbf{L}\) is summable in \(\mathbf{G}'\) .
\end{enumerate}

For a filter on G to be a Cauchy filter (resp. a convergent filter), it is necessary and sufficient that its image under the canonical mapping \(G \rightarrow G'\) be a Cauchy (resp. convergent) filter (General Topology, Chapter II, §3, no. 1, Proposition 4); whence first of all the equivalence of (a) and (b); on the other hand, as \(G'\) is metrizable, the equivalence of (b) and (c) follows from Proposition 9 of General Topology, Chapter IX, §2, no. 6.

Suppose now that \(G\) is commutative. Suppose that \(G'\) is complete and let \((x_{\lambda})_{\lambda \in L}\) be a family of elements of \(G'\) which converge to 0 with respect to \(\mathfrak{F}\) . For every neighbourhood \(V'\) of 0 in \(G'\) which is a subgroup of \(G'\) , there exists a finite subset \(J\) of \(L\) such that the condition \(\lambda \in L - J\) implies \(x_{\lambda} \in V'\) ; then \(\sum_{\lambda \in H} x_{\lambda} \in V'\) for every finite subset \(H\) of \(L\) not meeting \(J\) , which shows that the family \((x_{\lambda})_{\lambda \in L}\) is summable (General Topology, Chapter III, § 5, no. 2, Theorem 1).

Conversely, suppose that condition (d) holds and let \((x_{n})\) be a Cauchy sequence on \(G'\) ; the family \((x_{n+1} - x_{n})\) is then summable and in particular the series with general term \(x_{n+1} - x_{n}\) is convergent and hence the sequence \((x_{n})\) is convergent.

Let G be a filtered group whose filtration \((\mathbf{G}_{n})\) consists of normal subgroups of G; the quotient groups \(G/G_{n}\) are discrete and hence complete, since the \(G_{n}\) are open in G. Let \(f_{n}\) be the canonical mapping \(G \to G/G_{n}\) and for \(m \leqslant n\) let \(f_{mn}\) be the canonical mapping \(G/G_{n} \to G/G_{m}\) ; \((\mathbf{G}/\mathbf{G}_{n}, f_{mn})\) is an inverse system of discrete groups with Z as indexing set (General Topology, Chapter III, §7, no. 3). Let \(\tilde{G}\) be the topological group the inverse limit of this inverse system and for all n let \(g_{n}: \tilde{G} \to G/G_{n}\) be the canonical mapping; let \(f: G \to \tilde{G}\) be the inverse limit of the inverse system of mappings \((f_{n})\) such that \(f_{n} = g_{n} \circ f\) for all n; finally, let j be the canonical mapping of G to its Hausdorff completion \(\hat{G}\) ; as the \(G/G_{n}\) are complete, there exists a unique topological group isomorphism i: \(\hat{G} \to \tilde{G}\) such that \(f = i \circ j\) (loc. cit., Corollary 1 to Proposition 2); we shall call it the canonical isomorphism of \(\hat{G}\) onto \(\tilde{G}\) .

Let H be another filtered group whose filtration \((\mathbf{H}_{n})\) consists of normal subgroups of H and let \(u: G \to H\) be a homomorphism compatible with the filtrations (no. 4). Set \(\tilde{H} = \lim H/H_{n}\) ; for all n, u defines by taking quotients a homomorphism \(u_{n}: G/G_{n} \xleftarrow{\leftarrow} H/H_{n}\) and the \(u_{n}\) obviously form an inverse system of mappings; set \(\tilde{u} = \lim u_{n}\) . Moreover let \(\tilde{H}\) be the Hausdorff completion of H and \(\hat{u}: \hat{G} \to \hat{H}\) the homomorphism derived from u by passing to the Hausdorff completions (General Topology, Chapter II, § 3, no. 7, Proposition 15). It follows immediately from the definitions that if \(\tilde{G}\) is identified with \(\tilde{G}\) and \(\tilde{H}\) with \(\tilde{H}\) by means of the canonical isomorphisms, \(\hat{u}\) is identified with \(\tilde{u}\) . We conclude in particular that, if, for all n, \(u_{n}\) is an isomorphism, then \(\hat{u}\) is an isomorphism of topological groups.

\textbf{Examples of complete filtered groups and rings}\par

\begin{enumerate}
\def\labelenumi{(\arabic{enumi})}
\item
  Let G be a complete filtered group. Every closed subgroup of G with the induced filtration is complete (General Topology, Chapter II, § 3, no. 4, Proposition 8). Every quotient group of G with the quotient filtration is complete (General Topology, Chapter IX, § 3, no. 1, Remark 1).
\item
  Let A be a filtered commutative ring whose filtration we denote by \((a_n)_{n\in \mathbf{Z}}\) ; let \(A'\) be the ring of formal power series \(A[[X_1,\ldots ,X_s]]\) . For all \(e = (e_1,\dots ,e_s)\in \mathbf{N}^s\) , we write \(|e| = \sum_{i\in 1}^{s}e_i\) , \(X^e = \prod_{i\in 1}^{s}X_i^{e_i}\) so that every element \(P\in A'\) can be written uniquely \(P = \sum_{e\in \mathbf{N}^s}\alpha_{e,P}X^e\) where \(\alpha_{e,P}\in A\) . For all \(n\in \mathbf{Z}\) , let \(\mathfrak{a}_n'\) denote the set of \(\mathrm{P} \in \mathrm{A}'\) such that \(\alpha_{e,\mathrm{P}} \in \mathfrak{a}_{n-|e|}\) for all \(e \in \mathbf{N}^s\) ; we show that \(\mathfrak{a}_n'\) is an ideal of \(\mathrm{A}'\) . Clearly \(\mathfrak{a}_n'\) is an additive subgroup of \(\mathrm{A}'\) ; on the other hand, if \(\mathrm{P} \in \mathfrak{a}_n'\) and \(\mathrm{Q} \in \mathrm{A}'\) , then, for all \(e \in \mathbf{N}^s\) , \(\alpha_{e,\mathrm{PQ}} = \sum_{e' + e'' = e} \alpha_{e',\mathrm{Q}} \alpha_{e'',\mathrm{P}}\) ; as the relation \(e' + e'' = e\) implies \(|e''| \leqslant |e|\) , \(\mathrm{PQ} \in \mathfrak{a}_n'\) . Moreover, if \(\mathrm{Q} \in \mathfrak{a}_m'\) , then, for \(e' + e'' = e\) , \(\alpha_{e',\mathrm{Q}} \alpha_{e'',\mathrm{P}} \in \mathfrak{a}_{m-|e'|} \mathfrak{a}_{n-|e''}| \subset \mathfrak{a}_{m+n-|e|}\) , which proves that \((\mathfrak{a}_n')_{n \in \mathbf{Z}}\) is a filtration compatible with the ring structure on \(\mathrm{A}'\) (for obviously \(1 \in \mathfrak{a}_0'\) ). When in future we speak of \(\mathrm{A}'\) as a filtered ring, we shall mean, unless otherwise stated, with the filtration \((\mathfrak{a}_n')\) . Clearly \(\bigcap_{n \in \mathbf{Z}} \mathfrak{a}_n'\) is the set of formal power series all of whose coefficients belong to \(\bigcap_{n \in \mathbf{Z}} \mathfrak{a}_n\) ; then, if A is Hausdorff, so is \(\mathrm{A}'\) . If \(\mathfrak{a}_0 = \mathrm{A}\) , then \(\mathfrak{a}_0' = \mathrm{A}'\) .
\end{enumerate}

PROPOSITION 6. With the above notation, suppose that \(\mathfrak{a}_0 = \mathbf{A}\) and let \(h\) denote the mapping \(\mathbf{P} \mapsto (\alpha_{e,\mathbf{P}})_{e \in \mathbb{N}^s}\) . Then \(h\) is an isomorphism of the additive topological group \(\mathbf{A}'\) onto the additive topological group \(\mathbf{A}^{\mathbf{N}^s}\) . The polynomial ring \(\mathbf{A}[\mathbf{X}_1, \ldots, \mathbf{X}_s]\) is dense in \(\mathbf{A}'\) ; if \(\mathbf{A}\) is complete, so is \(\mathbf{A}'\) .

Clearly \(h\) is bijective; \(V_n = h(a_n')\) is the set of \((a_e) \in A^{N^s}\) such that \(a_e \in a_{n-|e|}\) for all \(e \in N^s\) such that \(|e| \leqslant n\) ; as these elements \(e\) are finite in number, \(V_n\) is a neighbourhood of 0 in \(A^{N^s}\) . Conversely, if \(V\) is a neighbourhood of 0 in \(A^{N^s}\) , there is a finite subset \(E\) of \(N^s\) and an integer \(v\) such that the conditions \(a_e \in a_v\) for all \(e \in E\) imply \((a_e) \in V\) ; if then \(n\) is the greatest of the integers \(v + |e|\) for \(e \in E\) , then \(h(a_n') \subset V\) , which proves the first assertion of Proposition 6. Moreover, with \(n\) and \(E\) defined as above, \(h(P - \sum_{e \in E} \alpha_{e, P} X^e) \in V\) for all \(P \in A'\) , which shows that \(A[X_1, \ldots, X_s]\) is dense in \(A'\) . The last assertion follows from the first and the fact that a product of complete spaces is complete.

Let \(m\) be an ideal of \(A\) and suppose that \((\mathfrak{a}_n)\) is the \(m\) -adic filtration; then, if \(n\) is the ideal of \(A'\) generated by \(m\) and the \(X_i\) ( \(1 \leqslant i \leqslant s\) ), the filtration \((\mathfrak{a}_n')\) is the \(n\) -adic filtration. For clearly, for all \(k \geqslant 0\) , \(n^k\) is generated by the elements \(aX^e\) such that \(a \in m^{k-|e|}\) for all \(e \in N^s\) such that \(|e| \leqslant k\) , whence \(n^k \subset \mathfrak{a}_k'\) . Let us prove conversely that \(\mathfrak{a}_k' \subset n^k\) . For all \(P \in \mathfrak{a}_k'\) , \(P = P' + P''\) , where \(P' = \sum_{|e| < k} \alpha_{e,P} X^e\) , \(P'' = \sum_{|e| \geqslant k} \alpha_{e,P} X^e\) . Clearly it is possible to write \(P'' = \sum_{|e|=k} X^e Q_e\) , where the \(Q_e\) are elements of \(A'\) , whence \(P'' \in n^k\) ; on the other hand, clearly \(\alpha_{e,P} X^e \in n^k\) for all \(e \in N^s\) , whence \(P' \in n^k\) . Then \(n^k = \mathfrak{a}_k'\) .

COROLLARY. Let A be a commutative ring,

\[
\mathbf {A} ^ {\prime} = \mathbf {A} [ [ \mathbf {X} _ {1}, \dots , \mathbf {X} _ {s} ] ]
\]

the ring of formal power series in s indeterminates over A and n the ideal of \(\Lambda'\) consisting of the formal power series with no constant term. The ring \(\Lambda'\) is Hausdorff and complete with the n-adic topology and the polynomial ring \(\mathbf{A}[\mathbf{X}_1,\dots ,\mathbf{X}_s]\) is everywhere dense in \(\mathbf{A}'\) .

It is sufficient to apply what has just been said to the case \(\mathfrak{m} = \{0\}\) .

